\newpage

\begin{figure}[h]
    \centering
    \includegraphics[width=1\linewidth]{4.png}
\end{figure}

Usamos la ecuación de las lentes delgadas:

\[
\frac{1}{f} = \frac{1}{s} + \frac{1}{s_i}
\]

Sustituyamos los valores:

\[
\frac{1}{-30} = \frac{1}{10} + \frac{1}{s_i}.
\]

Despejamos \( \frac{1}{s_i} \):

\[
\frac{1}{s_i} = -0.0333 - 0.1 = -0.1333.
\]

Invirtiendo...

\[
s_i = \frac{1}{-0.1333} \approx -7.5 \text{ cm}.
\]

Por lo tanto, la imagen se forma a 7.5 cm del centro óptico de la lente, en el mismo lado del objeto (lado izquierdo de la lente, lo que confirma que la imagen es virtual); es bueno usar las preguntas del ejercicio pasado para describir el problema así que las utilizare.

1. ¿Es real o virtual? 
   Como \( s_i \) es negativo, la imagen se encuentra en el mismo lado que el objeto. Esto indica que la imagen es virtual.

2. ¿Está invertida o derecha?  
   Las lentes divergentes siempre producen imágenes derechas. Nuevamente esto se puede ver con la magnificación

3. ¿Qué tan grande es respecto al objeto?  
   La magnificación \( M \) es dada como:

\[
   M = \frac{-(-7.5)}{10} = \frac{7.5}{10} = 0.75.
   \]

   La imagen es más pequeña que el objeto por un factor de 0.75 veces.  

4. ¿Cuál es su tamaño?  
   
   \[
   h' = 0.75 \times 6.00 = 4.50 \text{ cm}.
   \]

\begin{figure}[h]
    \centering
    \includegraphics[width=1\linewidth]{traz4.png}
\end{figure}
